\documentclass[a4paper,12pt]{article}
\usepackage[utf8]{inputenc}
\usepackage{amsmath}
\usepackage{graphicx}
\usepackage{listings}
\usepackage{xcolor}
\usepackage{geometry}
\geometry{margin=1in}

\definecolor{codegreen}{rgb}{0,0.6,0}
\definecolor{codegray}{rgb}{0.5,0.5,0.5}
\definecolor{codepurple}{rgb}{0.58,0,0.82}
\definecolor{backcolour}{rgb}{0.95,0.95,0.92}

\lstdefinestyle{verilogstyle}{
    backgroundcolor=\color{backcolour},   
    commentstyle=\color{codegreen},
    keywordstyle=\color{magenta},
    numberstyle=\tiny\color{codegray},
    stringstyle=\color{codepurple},
    basicstyle=\ttfamily\small,
    breakatwhitespace=false,         
    breaklines=true,                 
    captionpos=b,                    
    keepspaces=true,                 
    numbers=left,                    
    numbersep=5pt,                  
    showspaces=false,                
    showstringspaces=false,
    showtabs=false,                  
    tabsize=2,
    language=Verilog
}

\lstset{style=verilogstyle}

\title{Digital Logic Design: Project 2}
\author{Ashkan Namjoo}
\date{\today}


\begin{document}

\maketitle

\section{Project Github Repo:}
https://github.com/rowyfol/ds1-project-2026/tree/main/Pj2
\section{Introduction}
This report details the design and simulation of various digital components, culminating in the implementation of an 8-bit Ripple Carry Adder (RCA) and an 8-bit Carry Select Adder (CSA). All modules were described using Verilog with specific propagation delays to model realistic hardware behavior.

\section{Module Designs}

\subsection{Full Adder (1-bit)}
The foundation of our adders is the 1-bit Full Adder module (\texttt{full\_adder}). It takes two bits $A, B$ and a carry-in $C_{in}$ as inputs. We assigned a delay of $t_{sum} = 12$\,ns for the sum output and $t_{co} = 7$\,ns for the carry-out ($C_o$). These values are intentionally different from the original defaults and keep carry logic faster than sum logic, which is useful when comparing ripple-carry and carry-select timing paths.

\subsection{4-bit Ripple Carry Adder}
We designed a 4-bit adder (\texttt{adder4}) by cascading four 1-bit Full Adders. The carry output of each bit is linked to the carry input of the next bit. This module serves as a sub-component for our larger 8-bit designs.

\subsection{Multiplexers (1-bit and 4-bit)}
A 2-to-1 multiplexer (\texttt{mux2\_to\_1}) was implemented with a propagation delay of $t_{mux} = 5$\,ns. To handle multi-bit data, we created a 4-bit multiplexer (\texttt{mux4bit}) by instantiating four 1-bit MUXes in parallel.

\subsection{8-bit Adders}
Two different architectures were explored for 8-bit addition:
\begin{enumerate}
    \item \textbf{8-bit Ripple Carry Adder (RCA):} This is a straightforward extension of the ripple carry principle, using 8 Full Adder modules in a single chain.
    \item \textbf{8-bit Carry Select Adder (CSA):} This design uses one 4-bit adder for the lower bits and two separate 4-bit adders for the upper bits (pre-calculating results for both $C_{in}=0$ and $C_{in}=1$). A 4-bit MUX then selects the correct upper sum, and a 1-bit MUX selects the final $C_{out}$.
\end{enumerate}

\section{Simulation and Verification}
Each module was thoroughly tested to ensure logical correctness. We used testbenches that applied various input combinations, specifically focusing on three key scenarios:

\begin{itemize}
    \item \textbf{Scenario 1: Baseline Addition.} Adding small numbers (like $0x00 + 0x00$) to verify basic functionality.
    \item \textbf{Scenario 2: Carry Propagation.} Inputs like $0x0F + 0x01$ were used to ensure the carry correctly transitions from the lower nibble to the upper nibble.
    \item \textbf{Scenario 3: Maximum Load.} Adding $0xFF + 0xFF$ with an initial $C_{in}=1$ to test the circuit under maximum carry-out conditions.
\end{itemize}

The simulation results were captured in VCD (Value Change Dump) files and verified using GTKWave, confirming that all outputs matched the expected theoretical values.

\section{Delay Analysis and Comparison}

A critical part of this project was analyzing how different architectures affect performance. We used the following delay parameters: $t_{sum} = 12$\,ns, $t_{co} = 7$\,ns, and $t_{mux} = 5$\,ns.

\paragraph{Ripple Carry Adder (RCA):}
In an 8-bit RCA, the critical path is the carry propagation. The $n$-th sum bit is ready at $(n-1)t_{co} + t_{sum}$. For 8 bits ($n=8$):
\[ \text{Delay}_{RCA} = 7 \times 7\,\text{ns} + 12\,\text{ns} = 61\,\text{ns} \]
The final carry-out is ready at $8 \times 7\,\text{ns} = 56$\,ns.

\paragraph{Carry Select Adder (CSA):}
The CSA speeds things up by pre-calculating the upper 4 bits. The delay is determined by the lower 4-bit carry-out plus the multiplexer selection time:
\[ \text{Delay}_{CSA} = (4 \times 7\,\text{ns}) + 5\,\text{ns} = 33\,\text{ns} \text{ (for Cout)} \]
The upper sum bits are ready slightly later because the precomputed upper 4-bit sums require $3 \times 7\,\text{ns} + 12\,\text{ns} = 33$\,ns before the MUX selection, so their selected value is ready at $33\,\text{ns} + 5\,\text{ns} = 38$\,ns.

\paragraph{Final Verdict:}
The Carry Select Adder achieves a significant performance boost. By "spending" more hardware to pre-calculate potential outcomes, we reduced the total critical path delay from \textbf{61\,ns} to \textbf{38\,ns} (about a 38\% improvement).

\begin{table}[h]
\centering
\begin{tabular}{|l|c|c|}
\hline
\textbf{Metric} & \textbf{Ripple Carry Adder} & \textbf{Carry Select Adder} \\ \hline
Critical Path (Worst Case) & 61\,ns & 38\,ns \\ \hline
Complexity & Low (8 FAs) & High (12 FAs + MUXes) \\ \hline
Performance & Standard & Optimized \\ \hline
\end{tabular}
\caption{Comparison of 8-bit Adder Architectures}
\end{table}

This makes the CSA much more suitable for high-speed arithmetic units where performance is more critical than silicon area.

\section{Summary of Files}
All source files (\texttt{.v}) and testbenches are provided in their respective directories. The simulation waveforms for all 7 parts are included in the generated \texttt{.vcd} files, which can be viewed for detailed timing analysis.

\end{document}
